\subsection{正则局部环的商环}

命题 2. --- 设 A 是诺特局部环，\(\boldsymbol{x} = (x_1, ..., x_r)\) 是 \(\mathfrak{m}_{\mathrm{A}}\) 中元素的一个序列，且 \(\mathfrak{x}\) 是由 \(\boldsymbol{x}\) 生成的理想。下列条件等价：

\begin{enumerate}
\def\labelenumi{(\roman{enumi})}
\item
  A 是正则的，且 x 是 A 的一个坐标系的一部分；
\item
  环 A/x 是正则的，且 x 是 A 的一个正割序列；
\item
  环 A/x 是正则的，且 x 是 A 的一个完全正割序列。
\end{enumerate}

此外，在这些条件成立时，x 是 A 的一个素理想。

\begin{enumerate}
\def\labelenumi{(\roman{enumi})}
\setcounter{enumi}{2}
\item
  \(\Rightarrow\) (ii)：这由 § 3, no. 2 的命题 3 的推论得出。
\item
  \(\Rightarrow\) (i)：设 \(x\) 是 A 的一个正割序列，且诺特局部环 A/\(x\) 是正则的。设 \((x_{r+1}, ..., x_d)\) 是 A 中元素的一个序列，其模 \(x\) 的类构成 A/\(x\) 的一个坐标系。则序列 \((x_1, ..., x_d)\) 生成 A 的理想 \(m_A\)，且有
\end{enumerate}

\[
\dim (\mathrm{A}) = r + \dim (\mathrm{A} / \mathrm{x}) = r + (d - r) = d.
\]

因此，A 是正则的，且 \((x_{1},\ldots,x_{d})\) 是 A 的一个坐标系。

\begin{enumerate}
\def\labelenumi{(\roman{enumi})}
\tightlist
\item
  \(\Rightarrow\) (iii)：若条件 (i) 成立，则序列 \(x\) 是完全正割的（no. 2, 定理 1），从而根据 § 3, no. 2 的命题 3 的推论，它是正割的。因此有
\end{enumerate}

\[
\dim (\mathrm{A} / \mathfrak {x}) = \dim (\mathrm{A}) - r;\tag{3}
\]

此外，\(x_{1},\ldots,x_{r}\) 模 \(\mathfrak{m}_{\mathbf{A}}^{2}\) 的类在域 \(\kappa_{\mathbf{A}}\) 上线性无关，因此有

\[
\left[ \mathfrak {m} _ {\mathrm{A}} / (\mathfrak {m} _ {\mathrm{A}} ^ {2} + x): \kappa_ {\mathrm{A}} \right] = \left[ \mathfrak {m} _ {\mathrm{A}} / \mathfrak {m} _ {\mathrm{A}} ^ {2}: \kappa_ {\mathrm{A}} \right] - r.\tag{4}
\]

公式 (3) 和 (4) 表明 A/x 是正则的。

根据第 2 号的定理 1 的推论 1，每个正则诺特局部环都是整环。因此，如果 A/x 正则，则 x 是素理想。

推论 1. --- 设 A 是诺特局部环，且 t 是 \(\mathfrak{m}_{\mathrm{A}}\) 中的元素。下列条件等价：

\begin{enumerate}
\def\labelenumi{(\roman{enumi})}
\item
  A 是正则的，且 t 不属于 \(m_{A}^{2}\)；
\item
  A/tA 是正则的，且 \(\dim(\mathrm{A}/t\mathrm{A}) < \dim(\mathrm{A})\)；
\item
  A/tA 是正则的，且 t 不是 A 中的零因子。
\end{enumerate}

推论 2. --- 设 A 是正则诺特局部环，q 是 A 的一个理想。A/q 正则，当且仅当 q 由 A 的一个坐标系的某个子集生成。

充分性由命题 2 得出。

假设 A/q 正则，令 \(x = (x_{1}, \ldots, x_{r})\) 为 q 中元素的一个序列，其模 \(m_{A}^{2}\) 的类构成 \((\mathfrak{q} + \mathfrak{m}_{\mathrm{A}}^{2})/\mathfrak{m}_{\mathrm{A}}^{2}\) 关于域 \(\kappa_{A}\) 的一个基。令 x 为 A 中由 x 生成的理想。因此有 \(x \subset q\)，且 x 是 A 的一个坐标系的一部分，所以诺特局部环 A/x 是正则的（命题 2）；此外，向量空间 \(\mathfrak{m}_{\mathrm{A}}/(\mathfrak{q} + \mathfrak{m}_{\mathrm{A}}^{2})\) 与 \(\mathfrak{m}_{\mathrm{A}}/(\mathfrak{x} + \mathfrak{m}_{\mathrm{A}}^{2})\) 关于 \(\kappa_{A}\) 具有相同维数。因此，正则诺特局部环 A/q 和 A/x 具有相同维数。由于理想 q 和 x 都是素理想，且有 \(x \subset q\)，最后得到 q = x。

例。--- 设 \(k\) 为域，\(\mathbf{A} = k[[\mathbf{X}_1, ..., \mathbf{X}_n]]\)，\(\mathfrak{q}\) 是 \(\mathbf{A}\) 的理想且不同于 \(\mathbf{A}\)。\(\mathbf{A}/\mathfrak{q}\) 正则的充要条件是：存在整数 \(r \geqslant 0\) 以及元素 \(\mathbf{F}_1, ..., \mathbf{F}_r\)，这些元素属于 \(\mathbf{A}\)，生成 \(\mathfrak{q}\)，并且矩阵 \(\left(\frac{\partial\mathbf{F}_i}{\partial\mathbf{X}_j}(0, ..., 0)\right)\) 的秩为 \(r\)（“雅可比判据”）。于是 \(\dim(\mathbf{A}/\mathfrak{q}) = n - r\)。

注。--- 设 A 是正则诺特局部环，且 \(q \subset m_{A}\) 是 A 的一个理想，使得 A/q 正则。设 \((x_{1}, \ldots, x_{r})\) 是 q 中元素的一个序列，其模 \(m_{A}^{2}\) 的类生成向量空间 \((\mathfrak{q} + \mathfrak{m}_{\mathrm{A}}^{2}) / \mathfrak{m}_{\mathrm{A}}^{2}\)（关于域 \(\kappa_{A}\)）。命题 2 的推论 2 的证明表明，A 的理想 q 由 \((x_{1}, \ldots, x_{r})\) 生成。

